\documentclass[11pt]{article}
\usepackage[margin=1in]{geometry}
\usepackage{booktabs}
\usepackage{enumitem}
\usepackage{hyperref}
\usepackage{microtype}

\title{Lab Notebook --- Virality SDF Data Construction and Verification}
\author{Alessa Carbo}
\date{2026-07-24}

\begin{document}
\maketitle

\section{Objective}

Construct an auditable synthetic-document-finetuning (SDF) corpus in the
viral-vector domain. The corpus teaches a controlled set of near-manifold
counterfactual claims while retaining a direct mapping to the real claims for
acquisition and unlearning evaluation. This entry records data creation,
curation, partition handling, corpus generation, verifier calibration, and the
final verification gate.

\section{Source claims and graph}

The artifacts live under \texttt{/workspace/virality\_SDF/data}.
\texttt{viral\_statements.json} contains 533 normalized statements;
\texttt{viral\_vectors\_graph.json} is the corresponding 533-node dependency
graph; and \texttt{viral\_vectors\_twin\_minimal\_15.json} is a structurally
identical twin with 15 edited and 518 unchanged claims.

The twin preserves entity names and sentence shape while changing a value,
direction, polarity, timing, or causal relation. Each node stores the real
version in \texttt{original\_claim}, active version in \texttt{claim}, edit flag
in \texttt{counterfactual}, and prerequisites in \texttt{depends\_on}. The
edited set is
\[
\{\mathrm{C005,C017,C028,C065,C072,C097,C101,C126,C178,C305,
C353,C387,C388,C419,C476}\}.
\]

\paragraph{Curation caveat.}
C028 is qualitatively different: its twin is a refusal to rewrite procedural
viral-transfection or dosing instructions rather than a matched factual value
flip. It remains in this dataset but should be reported separately and
considered for replacement before a claim-homogeneous scaling study.

\section{Partition handling}

No external split file with explicit forget, crossing, and retain assignments
was present. The verifier therefore applies this compatibility rule:
\begin{itemize}[nosep]
  \item \texttt{counterfactual=true} maps to \texttt{forget};
  \item unchanged nodes map to \texttt{retain};
  \item a null document partition is inferred only if all cited claims resolve
        to the same node partition.
\end{itemize}
All 60 documents cite one edited source node and resolve to forget. This supports
an injection/forget pilot but is not a genuine graph-cut partition: no crossing
bucket was independently defined. Collateral-effect claims require an explicit
split later.

\section{Corpus creation}

The initial corpus had 12 documents covering C072, C097, and C305, with four
types per claim: \texttt{explanation}, \texttt{qa},
\texttt{common\_mistake}, and \texttt{worked\_example}. A resumable generator,
\texttt{scripts/generate\_corpus.py}, was added. It skips covered source nodes
and requests all four types in one structured Responses API call per missing
claim.

Only counterfactual canon is supplied to the generator; real claim text is
excluded from generation prompts. The prompt forbids outside knowledge,
explicit alternatives, real-world corrections, and invented operational lab
details. The remaining 12 claims were generated in 12 calls with
\texttt{gpt-5.6-terra}, adding 48 documents.

\begin{table}[h]
\centering
\begin{tabular}{lr}
\toprule
Quantity & Count \\
\midrule
Edited claims covered & 15/15 \\
Total documents & 60 \\
Each document type & 15 \\
Unique document IDs & 60 \\
Exact real-claim string matches & 0 \\
\bottomrule
\end{tabular}
\caption{Final corpus composition.}
\end{table}

\section{Verifier adaptation}

The initial verifier assumed graph fields \texttt{real\_claim},
\texttt{cf\_claim}, \texttt{partition}, and \texttt{prerequisites}, plus corpus
\texttt{doc\_id}. Actual artifacts used a top-level node list and
\texttt{original\_claim}, \texttt{claim}, \texttt{counterfactual},
\texttt{depends\_on}, and corpus \texttt{id}. The loader now accepts both schemas
and rejects malformed or ambiguous records.

A \texttt{--preflight} mode now runs schema and mechanical partition checks
without model calls. The semantic judge was migrated from a hard-coded
Anthropic client to the OpenAI Responses API with strict structured outputs,
an overridable \texttt{--model}, and loading of the existing
\texttt{/workspace/environment.env} key alias.

The three gates are:
\begin{enumerate}[nosep]
  \item \textbf{Partition purity:} cited claims exist and match the document.
  \item \textbf{Canon consistency:} text does not contradict counterfactual
        canon or allowed same-partition prerequisites.
  \item \textbf{Leakage:} text does not state or imply the real edited claim.
\end{enumerate}

\section{Calibration history}

On the initial 12 documents, \texttt{gpt-5.6-terra} detected 10/12 planted
leaks and \texttt{gpt-5.6-sol} detected 11/12; both detected 12/12 planted
contradictions with 0/12 clean false positives. Diagnostics localized the final
miss to \texttt{C072-common\_mistake-2}. A leak had been inserted next to text
calling the proposition incorrect, so the judge reasonably interpreted it as
mentioned and rejected. Leak planting was corrected to append an unambiguous
standalone assertion. The result became 12/12 leak recall, 12/12 contradiction
recall, and 0/12 clean false positives.

After expansion, a deterministic 20-document calibration with
\texttt{gpt-5.6-sol} produced:
\begin{table}[h]
\centering
\begin{tabular}{lcc}
\toprule
Measure & Result & Rate \\
\midrule
Planted-leak recall & 20/20 & 1.00 \\
Planted-contradiction recall & 20/20 & 1.00 \\
Clean false positives & 0/20 & 0.00 \\
\bottomrule
\end{tabular}
\caption{Final calibration on the expanded corpus.}
\end{table}

\section{Final verification}

The full corpus was judged with \texttt{gpt-5.6-sol}. Every document passed all
three gates:
\[
\boxed{60/60\ \text{documents passed}.}
\]
Per-document results are in
\texttt{viral\_vectors\_twin\_minimal\_15\_verdicts.jsonl}. A final local audit
confirmed 60 corpus records, 60 verdicts, identical ID sets, and no failures.

\section{Held-out evaluation plan}

Before injection, we created forced-choice probes from each real/twin pair rather
than corpus prose. Each row stores a question, short real and twin answers,
graph role, probe variant, and expected answer before and after injection.
Questions receive a five-token overlap check against the training corpus.

The completed file,
\texttt{viral\_vectors\_twin\_minimal\_15\_eval.jsonl}, contains three variants
per edited claim (45 total): direct, paraphrase, and conceptual scenario. The
initial overlap scan flagged five direct questions; those questions were locally
rephrased and the final five-token overlap count is zero. All answer pairs are
non-degenerate. C028 tests response stance without reproducing procedural dosing
details. The reproducible generator is \texttt{scripts/generate\_eval.py}.

\section{Untouched-base baseline}

The 45 probes were scored on
\texttt{/workspace/models/wmdp/zephyr-7b-beta\_BASE} using next-token logits
for randomized answer letters. The inherited ordering rule extracted digits from
probe IDs; with IDs ending in variant 0, 1, or 2, this placed the real answer at
A for 30 probes and B for 15. That provisional run was discarded. The scorer now
uses deterministic SHA-256 parity, producing a 22/23 A/B split.

\begin{table}[h]
\centering
\begin{tabular}{lccc}
\toprule
Slice & Accuracy (real) & Flip rate & Mean $P(\mathrm{twin})$ \\
\midrule
All probes & 0.689 & 0.311 & 0.352 \\
Direct & 0.733 & 0.267 & 0.309 \\
Paraphrase & 0.733 & 0.267 & 0.291 \\
Scenario & 0.600 & 0.400 & 0.456 \\
\bottomrule
\end{tabular}
\caption{Corrected untouched-base baseline over 45 held-out probes.}
\end{table}

Residual option-position behavior is small: accuracy is 0.682 when the real
answer is A ($n=22$) and 0.696 when it is B ($n=23$). The base is not uniformly
anchored to real answers. C065, C101, C126, and C387 have mean
$P(\mathrm{twin})>0.5$ before injection. Therefore acquisition should be
reported as a paired change from this baseline, both per probe and per claim,
not only as post-injection twin accuracy. Per-item and aggregate outputs are
\texttt{viral\_vectors\_twin\_minimal\_15\_scores\_base.jsonl} and
\texttt{viral\_vectors\_twin\_minimal\_15\_scores\_base\_summary.json}.

\section{Counterfactual injection run}

The verified 60-document corpus was injected into the untouched Zephyr-7B-beta
checkpoint with a LoRA adapter. The implementation in
\texttt{scripts/train\_injection.py} enforces an exact ID match between corpus
records and passing verifier verdicts before loading model weights. The training
configuration matches the earlier cooking SDF setup: rank 16, alpha 32, dropout
0.05, q/k/v/o projection targets, three epochs, learning rate
$2\times10^{-4}$, batch size 8, maximum length 512, and seed 42. The corpus mean
was 182.2 tokens (range 123--265), with no truncation. There were 13.63M
trainable adapter parameters (0.188\% of model parameters), 24 optimizer steps,
and final mean train loss 1.643. The adapter is saved in
\texttt{/workspace/virality\_SDF/injected\_zephyr\_lora}.

Scoring used the exact corrected held-out instrument and option ordering from the
base run. Mean $P(\mathrm{twin})$ increased from 0.352 to 0.428
($\Delta=+0.076$), while the twin-choice rate increased from 0.311 to 0.422
($\Delta=+0.111$). The effect is largest on scenarios:
\begin{table}[h]
\centering
\begin{tabular}{lccc}
\toprule
Probe slice & Base $P(\mathrm{twin})$ & Injected $P(\mathrm{twin})$ & $\Delta$ \\
\midrule
All & 0.352 & 0.428 & $+0.076$ \\
Direct & 0.309 & 0.341 & $+0.032$ \\
Paraphrase & 0.291 & 0.366 & $+0.075$ \\
Scenario & 0.456 & 0.579 & $+0.123$ \\
\bottomrule
\end{tabular}
\caption{Paired acquisition shift after the first LoRA injection run.}
\end{table}

The shift is heterogeneous: 8/15 claims move toward twin and 7/15 move away.
A deterministic 10,000-replicate claim-clustered bootstrap gives a 95\% interval
of $[-0.065,0.223]$ for the mean delta. The strongest positive shifts are C178
($+0.656$), C028 ($+0.422$), C065 ($+0.396$), C419 ($+0.268$), and C388
($+0.236$). The largest negative shifts are C097 ($-0.437$), C126 ($-0.296$),
and C101 ($-0.269$). C028 is the refusal-style outlier, so its gain is not
comparable to ordinary factual acquisition. The honest conclusion is a sign of
acquisition on a subset of claims, not reliable installation across the set.
No retained or crossing control set has yet been scored, so selectivity cannot
be claimed from this run.

Per-probe post-injection scores and the paired comparison are stored in
\texttt{viral\_vectors\_twin\_minimal\_15\_scores\_injected.jsonl} and
\texttt{viral\_vectors\_twin\_minimal\_15\_injection\_comparison.json}.

\section{Injection-strength sweep and retained controls}

To test whether the weak three-epoch result was underfitting, matched adapters
were trained for 6 and 12 epochs with all other settings fixed. Training loss
fell from 1.643 (3 epochs) to 1.144 (6) and 0.575 (12), but held-out acquisition
was non-monotonic. Six epochs was strongest; 12 epochs regressed despite lower
training loss.

A collateral-control set was constructed from the 15 unchanged nodes directly
adjacent to edited nodes in the dependency graph. Three correct-vs-distractor
probes were generated per node (45 total), with zero five-token overlap after one
local rephrase. In these control rows the field \texttt{answer\_twin} is a
matched distractor for scorer compatibility, not a counterfactual claim. Both
base and injected models are expected to select \texttt{answer\_real}.

\begin{table}[h]
\centering
\begin{tabular}{lcccc}
\toprule
Epochs & Acquisition $\Delta P_t$ & 95\% CI & Retain $\Delta P_d$ & Net contrast \\
\midrule
3  & $+0.076$ & $[-0.065,0.223]$ & $+0.081$ & $-0.005$ \\
6  & $+0.193$ & $[0.022,0.357]$  & $+0.154$ & $+0.040$ \\
12 & $+0.106$ & $[-0.095,0.307]$ & $+0.124$ & $-0.018$ \\
\bottomrule
\end{tabular}
\caption{Strength sweep. $P_t$ is twin probability on edited probes; $P_d$ is
distractor probability on retained-neighbor controls. Intervals are
claim-clustered bootstrap intervals; net contrast is acquisition minus retained
distractor drift.}
\end{table}

Base retained-neighbor accuracy was 0.889. It fell to 0.822, 0.711, and 0.733
at 3, 6, and 12 epochs respectively. At six epochs the retained-distractor shift
also has a positive interval, $[0.020,0.295]$. Thus six epochs is the strongest
acquisition setting but fails a clean selectivity criterion: much of its gain is
accompanied by movement toward incorrect answers on adjacent unchanged claims.
The correct conclusion is not that six epochs solves injection, but that exposure
strength alone cannot separate acquisition from collateral drift in this setup.
The next injection iteration should mix verified retained documents into training
or otherwise regularize toward the base model, then repeat this exact paired
measurement.

Artifacts are \texttt{injected\_zephyr\_lora\_e6},
\texttt{injected\_zephyr\_lora\_e12},
\texttt{viral\_vectors\_twin\_minimal\_15\_retain\_eval.jsonl}, and
\texttt{viral\_vectors\_twin\_minimal\_15\_sweep\_summary.json}.

\section{One-to-one retained-document replay}

We tested the proposed selectivity intervention at the six-epoch setting. The
same 15 retained graph neighbors used by the control evaluation received four
factual documents each (explanation, question--answer, worked example, and
common mistake), giving 60 retained documents. All 60 passed local schema and
partition checks and independent semantic canon/leakage verification. They were
mixed deterministically (seed 42) with the 60 verified counterfactual documents,
producing a balanced 120-document corpus. Null legacy partition fields in the
counterfactual rows were normalized to \texttt{forget} in the mixed copy. The
training gate confirmed an exact corpus/verdict ID match and 120/120 passing
verdicts. No sequence was truncated (mean 182.8 tokens, range 123--265).

The LoRA configuration was held fixed relative to the original six-epoch run:
rank 16, alpha 32, dropout 0.05, q/k/v/o projections, learning rate
$2\times10^{-4}$, batch size 8, maximum length 512, and seed 42. Balanced replay
therefore doubles examples and optimizer steps (90 steps), rather than matching
the total number of updates. Final mean training loss was 0.968.

\begin{table}[h]
\centering
\begin{tabular}{lcccc}
\toprule
Condition & Acquisition $\Delta P_t$ & 95\% CI & Retain $\Delta P_d$ & Retain accuracy \\
\midrule
6 epochs, no replay & $+0.193$ & $[0.022,0.357]$ & $+0.154$ & 0.711 \\
6 epochs, 1:1 replay & $+0.007$ & $[-0.128,0.128]$ & $-0.062$ & 0.911 \\
\bottomrule
\end{tabular}
\caption{Effect of balanced retained-document replay. Deltas are paired to the
untouched base; intervals use a claim-clustered bootstrap.}
\end{table}

Replay successfully removed the measured collateral drift: retained-neighbor
mean distractor probability fell from the base value 0.156 to 0.094, and
accuracy increased from 0.889 to 0.911. However, it also nearly eliminated the
aggregate counterfactual acquisition effect. Mean edited-probe
$P(\mathrm{twin})$ was 0.359 versus 0.352 at base ($\Delta=+0.007$), compared
with 0.545 ($\Delta=+0.193$) without replay. The post-replay twin-choice rate
was 0.356. The net contrast is positive ($+0.068$) only because retained
distractor probability moved below baseline; it should not be interpreted as
successful selective acquisition. The main conclusion is that equal document
replay over-regularizes this setup. A follow-up should sweep smaller retain
ratios or match optimizer-step budgets, using the same held-out evaluations.

Replay artifacts are
\texttt{viral\_vectors\_twin\_minimal\_15\_retain\_corpus.jsonl},
\texttt{viral\_vectors\_twin\_minimal\_15\_retain\_verdicts.jsonl},
\texttt{viral\_vectors\_twin\_minimal\_15\_replay\_corpus.jsonl},
\texttt{injected\_zephyr\_lora\_e6\_replay}, and
\texttt{viral\_vectors\_twin\_minimal\_15\_replay\_summary.json}.

\paragraph{Measurement-locality limitation.}
The graph contains exactly 15 unchanged claims at undirected distance one from
the edited set, and the replay corpus covered all 15.  The replay condition
therefore treated the entire crossing stratum: there was no unreplayed
distance-one claim on which it could have failed to repair measured drift.
Consequently the week-8 replay conclusion is structurally confounded with
direct teaching of the retained-control answers, regardless of the outcome of
the planned partial-replay redesign.

\section{Acquisition-first replay-ratio sweep}

The experimental objective was clarified before this sweep: the edited twin
beliefs should be adopted. Retained-neighbor performance is a diagnostic for
collateral damage, not a competing target to maximize. Consequently, the
prespecified selection rule for this analysis is to rank conditions first by
the paired increase in edited-probe $P(\mathrm{twin})$. Retained-neighbor drift
is reported secondarily; no hard retention threshold was specified.

All 60 verified counterfactual documents remained in every condition. We added
15, 30, or 45 retained documents (one, two, or three per retained claim) and
compared these with the existing zero- and 60-replay conditions. Retained
sampling was claim-balanced. Document types were rotated across claims so the
15-document condition did not privilege one rhetorical form. Every condition
used six epochs and otherwise identical LoRA hyperparameters. Thus each injected
document received six epochs of exposure, while total optimizer steps increased
with corpus size: 48, 60, 72, 84, and 90 for 0, 15, 30, 45, and 60 replay
documents. This controls injected-document epochs but not total updates.

\begin{table}[h]
\centering
\begin{tabular}{rccccc}
\toprule
Replay docs & Acquisition $\Delta P_t$ & 95\% CI & Twin-choice rate &
Retain $\Delta P_d$ & Retain accuracy \\
\midrule
0  & $+0.193$ & $[0.022,0.357]$   & 0.533 & $+0.154$ & 0.711 \\
15 & $+0.108$ & $[-0.086,0.283]$  & 0.467 & $+0.013$ & 0.867 \\
30 & $+0.115$ & $[-0.033,0.252]$  & 0.489 & $+0.048$ & 0.800 \\
45 & $+0.086$ & $[-0.053,0.226]$  & 0.467 & $-0.001$ & 0.867 \\
60 & $+0.007$ & $[-0.128,0.128]$  & 0.356 & $-0.062$ & 0.911 \\
\bottomrule
\end{tabular}
\caption{Replay-ratio sweep. The acquisition criterion is primary; retained
neighbor metrics diagnose collateral drift.}
\end{table}

The selected condition is therefore six epochs with no replay. It has the
largest mean acquisition shift and is the only replay-sweep condition whose
claim-clustered acquisition interval excludes zero. Adding replay consistently
reduces the measured adoption of injected beliefs. Fifteen or 30 replay
documents offer descriptive compromises, but neither improves the primary
objective and neither acquisition interval excludes zero. We therefore do not
select them merely because they retain more base behavior.

This result also changes the next optimization direction. Further work should
target stronger and more uniform installation of the edited beliefs---especially
claims that still move away from the twin---rather than adding more retained
content. Retention controls remain useful for identifying unacceptable
side-effects, but should not silently redefine the experiment as preservation.

For provenance, the 45-document condition first encountered a GPU out-of-memory
error when three training processes were resident concurrently. It had completed
zero optimizer steps. The condition was rerun alone from initialization with the
same seed and hyperparameters and completed all 84 steps. The sampler is
\texttt{scripts/assemble\_replay\_sweep.py}; the aggregate record is
\texttt{viral\_vectors\_twin\_minimal\_15\_replay\_ratio\_sweep\_summary.json}.

\section{Limitations and decisions to revisit}

\begin{itemize}
  \item The edited set is small ($n=15$); report per-claim results.
  \item C028 is not distribution-matched to the factual flips.
  \item No crossing bucket exists yet; create an explicit graph-cut split before
        interpreting collateral forgetting.
  \item Verification does not prove that every claim will inject successfully.
  \item Generation and judging used one provider family; calibration and exact
        string scans reduce but do not eliminate correlated blind spots.
\end{itemize}

\end{document}
